\section{Fundamentos de Q-Learning}

\subsection{Función Q óptima}

\begin{importantbox}{Ecuación óptima de Bellman}
La función Q óptima $Q^*$ satisface:
$$Q^*(s, a) = r(s, a) + \gamma \max_{a'} Q^*(s', a')$$

Propiedad clave: la política $\pi^*(s) = \arg\max_a Q^*(s, a)$ correspondiente a $Q^*$ es óptima entre todas las políticas.
\end{importantbox}

A diferencia de actor-critic, donde se ajusta $Q^\pi$ (la función Q de la política actual), Q-Learning ajusta directamente $Q^*$ (la función Q de la política óptima), usando el máximo en lugar de $\mathbb{E}_{a' \sim \pi}$.

\subsection{Algoritmo DQN}

Deep Q-Network (DQN) usa una red neuronal profunda para parametrizar $Q_\phi(s, a)$:

$$\min_\phi \sum_{(s,a,r,s') \in \mathcal{B}} \left\| Q_\phi(s, a) - \left(r + \gamma \max_{a'} Q_{\phi'}(s', a')\right) \right\|^2$$

\begin{importantbox}{Técnicas clave de DQN}
\begin{itemize}
    \item \textbf{Replay Buffer}: almacena todas las transiciones históricas y muestrea aleatoriamente para romper la correlación entre los datos.
    \item \textbf{Target Network}: usa un $Q_{\phi'}$ de actualización retrasada para calcular el TD target, evitando la inestabilidad de «perseguirse a sí mismo».
    \item \textbf{Exploración $\epsilon$-Greedy}: con probabilidad $\epsilon$ se elige una acción aleatoria, de lo contrario se elige $\arg\max_a Q(s,a)$.
\end{itemize}
\end{importantbox}

\subsection{Resumen de la sección}
 
Q-Learning obtiene la política óptima de manera indirecta al ajustar la función Q óptima. DQN extiende Q-Learning a espacios de estados de alta dimensión; replay buffer y target network son las técnicas clave de estabilización.

